\chapter{Literature Review}
\label{chap:litreview}

\section[Scope and Criteria]{Scope and Criteria}

The broader project surrounding this thesis contains three connected stages: risk quantification, asset-universe selection, and final weight allocation. The present thesis covers the middle stage only. Accordingly, this review does not treat portfolio weighting, rebalancing, or risk-scoring frameworks as primary literature targets unless they help clarify the boundary of the selection problem. The thesis consumes quantified risk information from an earlier stage and then uses investor-profile information to support asset-universe selection before final allocation. The central question is therefore narrower: how can an AI-supported system choose a suitable investable universe for an investor whose risk tolerance is already given as an input and whose market risk context has already been quantified upstream?

This scope changes the selection logic of the reviewed papers. Priority is given to studies that perform screening, ranking, filtering, or preselection before a later optimization stage. A second priority is given to studies in which investor profile or risk appetite explicitly changes the recommendation outcome. A third priority is given to AI and machine-learning methods, because the thesis itself is framed as an AI-based selection problem. Under these criteria, the literature is best read through three themes: preselection as a distinct financial decision stage, investor-conditioned automated recommendation, and the small set of studies that come closest to combining both ideas within a single pipeline.

\section[Preselection Before Allocation]{Preselection Before Allocation}

Long before machine-learning methods were used for asset selection, finance research had already treated candidate-universe construction as a distinct step. Factor-based portfolio traditions sorted or screened securities using common exposures such as size and value, creating structured subsets before later portfolio formation \cite{fama1993common}. Multi-criteria decision frameworks were likewise used to rank or filter candidate stocks on profitability, liquidity, and risk before final portfolio design \cite{vasquez2022ahptopsis}. These earlier streams matter here because they show that staged universe construction is not an AI artifact; the machine-learning literature reviewed below should be understood as extending and automating an already recognized financial decision stage.

Several studies support the claim that universe selection should be modeled separately from final weight allocation. A deep-learning portfolio-formation framework showed that stock preselection can be performed before downstream optimization, rather than being treated as an incidental preprocessing step \cite{wang2020portfolioformation}. A related study used random-forest return predictions to select a subset of stocks before applying mean-variance and hierarchical-risk-parity allocation rules, and the reported results indicated that this staged procedure improved the quality of the final portfolio relative to naive weighting rules \cite{kaczmarek2022building}. Another paper reached a similar conclusion from a broader model-comparison perspective by combining return prediction with portfolio formation and explicitly applying machine-learning and deep-learning models to stock preselection before optimization \cite{ma2021returnprediction}. Together, these papers establish that preselection is not merely a data-cleaning convenience. It is a financially meaningful modeling stage that can materially affect downstream portfolio performance.

The same separation appears in work that is not purely machine-learning based. A two-stage robust portfolio-construction framework first screened candidate stocks and only then solved the final allocation problem on the qualified set \cite{mills2022twostage}. Although that study is not personalized and does not map investor profile to the first-stage filter, it still supports the architectural argument that candidate-universe qualification and final weighting are distinct tasks. A further analogy appears in a stock-recommendation system that used machine learning to score and rank stocks after aggregating signals from different investor groups \cite{qin2024followakoinvestor}. That paper is not personalized to the target investor, because the investor groups act as information sources rather than as the receiving client profile. Even so, it remains useful because it frames recommendation as a ranking-based construction of a candidate stock set rather than as immediate weight assignment. However, the studies in this stream remain largely indifferent to the characteristics of the investor who will eventually hold the resulting portfolio. They improve candidate-universe quality, but they do not condition the first-stage filter on investor risk tolerance or other profile-dependent constraints. That unresolved limitation motivates the next stream of work.

\section[Personalized Recommendation]{Personalized Recommendation}

A separate body of work demonstrates that investor-specific information can meaningfully alter automated investment recommendations. A case-based finance advisory system used investor-specific information to personalize finance advice and diversification strategies \cite{musto2015personalized}. An AI-based investment-type recommender likewise used investor and expert feedback to map investor characteristics to suggested investment types \cite{asemi2024investmenttype}. A more recent fund-recommendation study introduced dynamic utility learning and generated personalized fund lists from a predefined product pool while updating user preferences through interaction data \cite{wei2025fundrecommendation}. Across these studies, the common pattern is that personalization changes what is recommended, but the output remains advice, product-type classification, or a fund list rather than a dedicated pre-optimization universe that can be passed to a separate allocation stage.

Recent generative-AI studies reinforce this point at the portfolio-output level. One study compared ChatGPT with robo-advisors across three investor profiles and reported that, under the benchmark and profile definitions used in that study, AI-generated recommendations aligned more closely with the reference advice \cite{oehler2024chatgptrobo}. Another study prompted ChatGPT to build stock portfolios for different risk appetites and found that the resulting selections varied systematically with the stated appetite for risk \cite{schneider2025riskappetite}. Together, these papers show that investor-profile-conditioned AI recommendation is technically feasible. Their shared limitation is that they output full portfolio advice or directly selected stock portfolios, so the preselection and final-allocation stages remain collapsed rather than separated.

\section[Closest Overlap and Gap]{Closest Overlap and Remaining Gap}

The closest work identified in this review is a deep reinforcement learning framework that combines investor-specific risk categories, AI-based asset selection, and downstream portfolio construction within one pipeline \cite{orra2025deep}. At the time of writing, this study is available as an arXiv preprint rather than a peer-reviewed journal article, so its findings should be treated as preliminary rather than established. It remains relevant because it connects three ideas that are often separated in the literature: investor profile, AI-based selection, and sequential financial decision-making. It is therefore the closest available comparison point, while also illustrating how little peer-reviewed work has yet combined these elements in a single framework.

At the same time, the literature as a whole still leaves a clear gap. Preselection papers usually improve candidate-universe construction without conditioning on investor profile \cite{wang2020portfolioformation, kaczmarek2022building, ma2021returnprediction, mills2022twostage, qin2024followakoinvestor}. Personalized recommendation papers usually condition on investor profile or learned utility, but they output investment advice, fund lists, product types, or full portfolios rather than a dedicated pre-optimization universe \cite{musto2015personalized, asemi2024investmenttype, wei2025fundrecommendation, oehler2024chatgptrobo, schneider2025riskappetite}. Even the closest overlap still operates in a stock-level setting and integrates the later portfolio stage directly \cite{orra2025deep}. The defended contribution of this thesis is therefore not the isolated idea of personalization, nor the isolated idea of AI-based selection. Its contribution is the combination of profile-aware, AI-supported asset-universe selection after upstream risk quantification and before downstream final weight allocation, with the problem framed explicitly as the middle stage of a broader portfolio pipeline.
